We set out to build, for all seven machine types of the DCASE 2026 Task~2 development set, an anomalous sound detector trained on normal recordings alone, and to run the three experiments the problem statement asked for. What exists at the end is a complete and reproducible pipeline: an audit of 8\,400 recordings, a shared two-channel log-Mel front end, fourteen trained Autoencoders, normal-distribution statistics for Mahalanobis scoring in two feature spaces, 2\,800 scored test recordings, 42 calibrated thresholds, full metric tables, verification reports and a demonstration application. A final audit recomputed every reported number from the raw score files and found no discrepancy.

The results themselves are sober. Residual-space Mahalanobis scoring with the source-plus-target model gave the best average performance (AUC 0.573, pAUC 0.537), reconstruction MSE was a close second, and scoring in the latent space was not useful. The source domain is handled reasonably; the target domain is not, and the ten target recordings help only a little. Calibrating thresholds on training scores, which the protocol leaves as the only option, proved clearly optimistic.

If we were to continue, we would, in this order: run the cross-fitted calibration stage that is already implemented, so that operating points rest on independent normal scores; evaluate selective Mahalanobis scoring with per-domain statistics and per-domain thresholds; repeat the training over several seeds to learn how much of the differences between configurations is noise; replace the dense Autoencoder with a representation trained to be discriminative, for example through attribute classification or contrastive pre-training, which recent challenge entries have found more robust to domain shift; and make explicit use of the two channels, now that the pipeline preserves both, through inter-channel features rather than simple concatenation.
